\subsection{Graph Reduction Techniques}
\label{sec:relwork:reduction}
Each family below is fixed by the quantity its operator controls, and the methods inside a
family differ only in how that quantity is chosen.

\subsubsection{Graph Partitioning}
\label{sec:relwork:reduction:partitioning}
A $k$-way partition $\pi \colon \mathcal{V} \rightarrow \{1, \dots, k\}$ deletes every edge
whose endpoints land in different blocks, so a partitioner controls the cut and nothing
else. The reference objective is the balanced minimum cut,
\begin{equation}
    \min_{\pi}\ \bigl| \{\, (u, v) \in \mathcal{E} : \pi(u) \neq \pi(v) \,\} \bigr|
    \quad \text{subject to} \quad
    \max_{i}\ \bigl| \pi^{-1}(i) \bigr| \leq (1 + \varepsilon)
    \left\lceil \frac{|\mathcal{V}|}{k} \right\rceil,
    \label{eq:relwork:mincut}
\end{equation}
with imbalance tolerance $\varepsilon$. The problem is NP-hard, and METIS is the standard
approximation to it \citep{karypis1997metis}: the graph is coarsened by successive
matchings, the smallest graph is partitioned directly, and the assignment is projected back
with local refinement at each level. Every cut edge in~\eqref{eq:relwork:mincut} counts one,
whatever it connects. That indifference is the only term the domain-aware variant
of~\ref{sec:method:partition:spanmetis} changes.

Two further constructions are references rather than competitors. Hash partitioning draws
$\pi(v)$ uniformly and independently of $\mathcal{E}$, which is the default in distributed
GNN systems, where partition quality is traded against the cost of computing it
\citep{zheng2020distdgl}; each edge then survives with probability $1/k$, and a
structure-aware objective earns its cost only by retaining more than that fraction.
Ordering-based slicing sorts $\mathcal{V}$ by a scalar the graph already carries and cuts
the sorted sequence into $k$ equal blocks, so balance is exact and the cut follows entirely
from the ordering (\ref{sec:method:partition:levelslice}).

\subsubsection{Graph Sparsification}
\label{sec:relwork:reduction:sparsification}
Sparsification keeps a subset $\mathcal{E}' \subseteq \mathcal{E}$, or an induced subgraph
on a node subset $\mathcal{V}' \subseteq \mathcal{V}$, and the methods differ in what they
promise the survivor preserves. The strongest such promise available at near-linear cost is
a bound on stretch. A $t$-spanner is a subgraph $H \subseteq G$ with
\begin{equation}
    d_H(u, v) \leq t \cdot d_G(u, v)
    \qquad \text{for every } u, v \in \mathcal{V},
    \label{eq:relwork:stretch}
\end{equation}
so an edge may be discarded only when an alternative path of length at most $t$ survives
\citep{peleg1989spanners}, and randomised constructions reach the optimal size-stretch
trade-off in linear time \citep{baswana2007spanners}. The guarantee is paid for out of dense
short-range redundancy. A graph that lacks it returns a spanner nearly as large as itself,
which is the negative result reported in~\ref{sec:method:sparse:forest}.
% Cut: spectral sparsification makes the analogous promise about the Laplacian quadratic form.

The remaining methods trade the guarantee for a knob or for simplicity. A minimum-weight
spanning forest keeps $|\mathcal{V}| - c$ edges on $c$ components \citep{kruskal1956mst}: a
connectivity backbone with every cycle destroyed, at a compression that is a property of the
graph rather than a setting. Independent edge removal keeps each edge with probability
$1 - p$ and was introduced as a per-epoch regulariser for deep GNNs
\citep{rong2020dropedge}; a mask drawn once instead of per epoch is a reduction, and its
ratio is freely tunable, which is what lets it serve as a matched-compression control.
Centrality pruning keeps the $\lceil \kappa |\mathcal{V}| \rceil$ highest-scoring nodes
under PageRank \citep{brin1998anatomy,langville2003pagerank} and returns the induced
subgraph, so every edge incident on a discarded node goes with it, and $\kappa$ is again
free. The two ratios of~\eqref{eq:prelim:retention} separate at exactly this point: the
first two methods hold $\eta_{\mathcal{V}} = 1$, the last moves both.

\subsubsection{Graph Summarization and Coarsening}
\label{sec:relwork:reduction:summarization}
Coarsening merges nodes into super-nodes under the merge map
of~\ref{sec:prelim:reduction:summarization}, and the strands differ in what they treat as
grounds for merging two nodes.

The exact strand merges only nodes that are provably indistinguishable. Colour
refinement~\eqref{eq:prelim:cr} computes the coarsest equitable partition
\citep{grohe2014colourrefinement}, and \citet{bollen2023exact} turn that partition into a
training procedure: a graph is quotiented by its $d$-step refinement classes, and a network
trained on the quotient produces exactly the output it would have produced on the original.
Exactness is conditional, and the conditions are what make the result usable. The model must
be message passing of depth at most $d$ over an aggregator that reads the multiset rather
than the set of incoming messages, so that neighbour counts still reach it. The graph must
be presented as a multigraph whose super-edges carry those multiplicities, and its initial
features must be constant on each initial colour class. A count cap $c$ truncates the
multiset, with $c = 1$ collapsing it to a set and giving bisimulation: coarser, and no
longer exact for a counting model. Section~\ref{sec:prelim:reduction} states the condition
formally. Both requirements are met by construction here, since the encoder sums over
neighbours and the summarized schema records super-edge multiplicities
(\ref{sec:method:summary:schema}), so refinement at $c = \infty$ is lossless for this
encoder and is the study's positive control rather than one more contender
(\ref{sec:method:architecture:exact}). k-SNAP groups by attribute and neighbour class
\citep{tian2008ksnap} and is the depth-one case of the same construction, so it is cited and
not run.

The spectral and classical strand descends from scientific computing. Local Variation
coarsening contracts node sets under a restricted spectral approximation bound
\citep{loukas2019localvariation}, the lineage is surveyed by \citet{chen2022coarsening}, and
Kron reduction has a directed extension \citep{sugiyama2023kron}. Its guarantees are stated
for the spectrum of a symmetric Laplacian, so none of them survives the move to a directed
graph unexamined, and the strand is surveyed here rather than run.

Convolution matching is the GNN-aware strand \citep{dickens2024convmatch}: nodes merge when
the graph convolution itself cannot separate them, rather than when some graph property
cannot, and about 95\% of full-graph performance is reported at 1\% of the node count on
node classification. Merging under the operator the model actually applies is the strongest
general-purpose criterion published, which makes it the bar a domain-specific method has to
clear (\ref{sec:method:summary:convmatch}). Universal Graph Coarsening merges nodes that
collide under locality-sensitive hashing, in time linear in the size of the graph
\citep{ugc2024}, and occupies the cheap end of the same axis.
\citet{shabani2023summarization} survey the use of both strands with GNNs.

Uniform merging to a stated target is the control the rest are read against. Compression is
a free parameter for only some coarseners, so no general procedure makes two real methods
meet at one ratio; a random merge reaches any ratio exactly, and each method is therefore
compared against a random control run at that method's own achieved compression
(\ref{sec:method:summary:random}).

Condensation is the family this thesis excludes (\ref{sec:intro:scope:reductions}). GCond
and its descendants synthesise a new small graph by gradient matching instead of merging the
real one \citep{jin2022gcond}, so (i) no correspondence exists between synthetic and real
nodes, (ii) the method is label-dependent by construction, (iii) gradient matching ties the
result to one architecture, and (iv) with no real nodes there is nothing to run cross-state
inference on.

\subsubsection{Reduction as Preprocessing for GNN Training}
\label{sec:relwork:reduction:forgnn}
Two systems already train on a reduced graph and then serve the original. SCAL coarsens the
training graph and applies the same weights to the full graph at inference
\citep{huang2021scal}, the shared-weights mechanism RQ5 tests.
\citet{generale2022scaling} instead transfer weights back through the node-to-super-node
mapping for a node-level task on knowledge graphs, and their ablation finds the
\emph{content} carried on a super-node to be the deciding factor, which is why member counts
and level statistics travel with super-nodes here (\ref{sec:method:summary:mergemap}). The
distance between the two is the distance between their tasks. A node-level prediction needs
a mapping back to individual nodes, whereas a graph-level regression needs only an input
space that both structural states inhabit (\ref{sec:method:summary:schema}), so shared
weights suffice and no transfer step exists to go wrong.
